% GENERATED FILE. Edit canonical lessons, then run npm run generate:books.
% canonical-source-hash: fnv1a64:3c3dde3286e2fea9
% canonical-lessons: SA-C541-dolayate, SA-C541-vimrsati, SA-C541-dvidha, SA-C541-aparah, SA-C541-vibhinnah, SA-R541-review-534-541

\chapter{Two Ways to Go}
\label{ch:sa-two-ways-to-go}
\hlchaptermodality{541}

\begin{chapteropening}
\textbf{By the end of this chapter:} \emph{I can weigh two options and say what I would do.}

Complete the last lesson of chapter 541: I can weigh two options and say what I would do.
\end{chapteropening}

\section[\texorpdfstring{\sk{दोलायते} (dolāyate)}{दोलायते (dolāyate)}]{\sk{दोलायते} (dolāyate) — wavers}

\label{lesson:SA-C541-dolayate}



\begin{quote}
Before the new one: say the Sanskrit for a waste of time, then the Sanskrit for spends time.
\end{quote}

\subsection*{You'll want to know: \sk{दोलायते}}
\textbf{\sk{दोलायते}} — \emph{dolāyate} — \textquotedblleft{}wavers\textquotedblright{}.

\begin{quote}
\sk{यदि} \sk{मनः} \sk{दोलायते}, \sk{तर्हि} \sk{मित्रं} \sk{पृच्छतु।}
\end{quote}

\begin{quote}
\emph{yadi manaḥ dolāyate, tarhi mitraṃ pṛcchatu.}
\end{quote}

\begin{quote}
If your mind wavers, ask a friend.
\end{quote}

Past (he …): \sk{सः} \sk{बहु} \sk{दोलायितवान्}, \emph{saḥ bahu dolāyitavān}.

\begin{cousinweb}
\textbf{\sk{दोलायते}} (\emph{dolāyate}) is built on \emph{dolā}, \textquotedblleft{}a swing\textquotedblright{}: MW gives \emph{dolāyate} as to rock about like a swing, to move to and fro, to be doubtful (MW).
\end{cousinweb}

\begin{grammarlens}[title={What you've built}]
One more word for saying what happens if.
\end{grammarlens}

\subsection*{Your turn}
\begin{itemize}
\raggedright
  \item \emph{Say it:} \emph{dolāyate}
  \item \emph{Say it:} \emph{dolāyate}, once more
  \item \emph{Say it:} \emph{yadi manaḥ dolāyate, tarhi mitraṃ pṛcchatu.} — If your mind wavers, ask a friend.
  \item \emph{Recall:} say the Sanskrit for a waste of time, then the Sanskrit for spends time, then say \emph{dolāyate} again
\end{itemize}

\begin{tcolorbox}[breakable,colback=teal!4,colframe=teal!35!black,title={Before you move on}]
What does \sk{दोलायते} mean? (\textquotedblleft{}Wavers\textquotedblright{}.) Say it once more.
\end{tcolorbox}
\section[\texorpdfstring{\sk{विमृशति} (vimṛśati)}{विमृशति (vimṛśati)}]{\sk{विमृशति} (vimṛśati) — weighs up}

\label{lesson:SA-C541-vimrsati}



\begin{quote}
Before the new one: say the Sanskrit for spends time, then the Sanskrit for wavers.
\end{quote}

\subsection*{You'll want to know: \sk{विमृशति}}
\textbf{\sk{विमृशति}} — \emph{vimṛśati} — \textquotedblleft{}weighs up\textquotedblright{}.

\begin{quote}
\sk{यदि} \sk{त्वं} \sk{विमृशसि}, \sk{तर्हि} \sk{उचितः} \sk{मार्गः} \sk{स्पष्टः} \sk{भविष्यति।}
\end{quote}

\begin{quote}
\emph{yadi tvaṃ vimṛśasi, tarhi ucitaḥ mārgaḥ spaṣṭaḥ bhaviṣyati.}
\end{quote}

\begin{quote}
If you think it over, the right way will become clear.
\end{quote}

Past (he …): \sk{सः} \sk{चिरं} \sk{विमृष्टवान्}, \emph{saḥ ciraṃ vimṛṣṭavān}.

\begin{cousinweb}
\textbf{\sk{विमृशति}} (\emph{vimṛśati}) is built on \emph{vi}, \textquotedblleft{}apart\textquotedblright{}, and the root √mṛś, \textquotedblleft{}to touch\textquotedblright{}: to touch with the mind; MW gives \emph{vimṛś} as to consider, to reflect on, to deliberate; its past participle is \emph{vimṛṣṭa} (MW).
\end{cousinweb}

\begin{grammarlens}[title={What you've built}]
One more word for saying what happens if.
\end{grammarlens}

\subsection*{Your turn}
\begin{itemize}
\raggedright
  \item \emph{Say it:} \emph{vimṛśati}
  \item \emph{Say it:} \emph{vimṛśati}, once more
  \item \emph{Say it:} \emph{yadi tvaṃ vimṛśasi, tarhi ucitaḥ mārgaḥ spaṣṭaḥ bhaviṣyati.} — If you think it over, the right way will become clear.
  \item \emph{Recall:} say the Sanskrit for spends time, then the Sanskrit for wavers, then say \emph{vimṛśati} again
\end{itemize}

\begin{tcolorbox}[breakable,colback=teal!4,colframe=teal!35!black,title={Before you move on}]
What does \sk{विमृशति} mean? (\textquotedblleft{}Weighs up\textquotedblright{}.) Say it once more.
\end{tcolorbox}
\section[\texorpdfstring{\sk{द्विधा} (dvidhā)}{द्विधा (dvidhā)}]{\sk{द्विधा} (dvidhā) — in two ways}

\label{lesson:SA-C541-dvidha}



\begin{quote}
Before the new one: say the Sanskrit for wavers, then the Sanskrit for weighs up.
\end{quote}

\subsection*{You'll want to know: \sk{द्विधा}}
\textbf{\sk{द्विधा}} — \emph{dvidhā} — \textquotedblleft{}in two ways\textquotedblright{}.

\begin{quote}
\sk{यदि} \sk{मार्गः} \sk{द्विधा} \sk{भवति}, \sk{तर्हि} \sk{कुत्र} \sk{गच्छामि}?
\end{quote}

\begin{quote}
\emph{yadi mārgaḥ dvidhā bhavati, tarhi kutra gacchāmi?}
\end{quote}

\begin{quote}
If the road splits in two, which way do I go?
\end{quote}

\begin{cousinweb}
\textbf{\sk{द्विधा}} (\emph{dvidhā}) is built on \emph{dvi}, \textquotedblleft{}two\textquotedblright{}, with the ending \emph{-dhā}, \textquotedblleft{}in ... ways\textquotedblright{}; MW gives \emph{dvidhā} as in two ways or parts, divided (MW).
\end{cousinweb}

\begin{grammarlens}[title={What you've built}]
One more word for saying what happens if.
\end{grammarlens}

\subsection*{Your turn}
\begin{itemize}
\raggedright
  \item \emph{Say it:} \emph{dvidhā}
  \item \emph{Say it:} \emph{dvidhā}, once more
  \item \emph{Say it:} \emph{yadi mārgaḥ dvidhā bhavati, tarhi kutra gacchāmi?} — If the road splits in two, which way do I go?
  \item \emph{Recall:} say the Sanskrit for wavers, then the Sanskrit for weighs up, then say \emph{dvidhā} again
\end{itemize}

\begin{tcolorbox}[breakable,colback=teal!4,colframe=teal!35!black,title={Before you move on}]
What does \sk{द्विधा} mean? (\textquotedblleft{}In two ways\textquotedblright{}.) Say it once more.
\end{tcolorbox}
\section[\texorpdfstring{\sk{अपरः} (aparaḥ)}{अपरः (aparaḥ)}]{\sk{अपरः} (aparaḥ) — the other one}

\label{lesson:SA-C541-aparah}



\begin{quote}
Before the new one: say the Sanskrit for weighs up, then the Sanskrit for in two ways.
\end{quote}

\subsection*{You'll want to know: \sk{अपरः}}
\textbf{\sk{अपरः}} — \emph{aparaḥ} — \textquotedblleft{}the other one\textquotedblright{}.

\begin{quote}
\sk{यदि} \sk{एषः} \sk{मार्गः} \sk{कठिनः} \sk{अस्ति}, \sk{तर्हि} \sk{अपरं} \sk{मार्गं} \sk{गच्छामः।}
\end{quote}

\begin{quote}
\emph{yadi eṣaḥ mārgaḥ kaṭhinaḥ asti, tarhi aparaṃ mārgaṃ gacchāmaḥ.}
\end{quote}

\begin{quote}
If this road is hard, we take the other one.
\end{quote}

\begin{cousinweb}
\textbf{\sk{अपरः}} (\emph{aparaḥ}) comes from \emph{apa}, \textquotedblleft{}away, off\textquotedblright{}, as MW gives it; MW records \emph{apara} as later, latter, the other (MW).
\end{cousinweb}

\begin{grammarlens}[title={What you've built}]
One more word for saying what happens if.
\end{grammarlens}

\subsection*{Your turn}
\begin{itemize}
\raggedright
  \item \emph{Say it:} \emph{aparaḥ}
  \item \emph{Say it:} \emph{aparaḥ}, once more
  \item \emph{Say it:} \emph{yadi eṣaḥ mārgaḥ kaṭhinaḥ asti, tarhi aparaṃ mārgaṃ gacchāmaḥ.} — If this road is hard, we take the other one.
  \item \emph{Recall:} say the Sanskrit for weighs up, then the Sanskrit for in two ways, then say \emph{aparaḥ} again
\end{itemize}

\begin{tcolorbox}[breakable,colback=teal!4,colframe=teal!35!black,title={Before you move on}]
What does \sk{अपरः} mean? (\textquotedblleft{}The other one\textquotedblright{}.) Say it once more.
\end{tcolorbox}
\section[\texorpdfstring{\sk{विभिन्नः} (vibhinnaḥ)}{विभिन्नः (vibhinnaḥ)}]{\sk{विभिन्नः} (vibhinnaḥ) — different}

\label{lesson:SA-C541-vibhinnah}



\begin{quote}
Before the new one: say the Sanskrit for in two ways, then the Sanskrit for the other one.
\end{quote}

\subsection*{You'll want to know: \sk{विभिन्नः}}
\textbf{\sk{विभिन्नः}} — \emph{vibhinnaḥ} — \textquotedblleft{}different\textquotedblright{}.

\begin{quote}
\sk{यदि} \sk{मतानि} \sk{विभिन्नानि} \sk{सन्ति}, \sk{तर्हि} \sk{वयं} \sk{पुनः} \sk{चिन्तयामः।}
\end{quote}

\begin{quote}
\emph{yadi matāni vibhinnāni santi, tarhi vayaṃ punaḥ cintayāmaḥ.}
\end{quote}

\begin{quote}
If our opinions differ, we think again.
\end{quote}

\begin{cousinweb}
\textbf{\sk{विभिन्नः}} (\emph{vibhinnaḥ}) is built on \emph{vi}, \textquotedblleft{}apart\textquotedblright{}, and the root √bhid, \textquotedblleft{}to split\textquotedblright{}, which also gives \emph{bhinatti}, as its past participle; MW gives \emph{vibhinna} as split, divided, and various (MW).
\end{cousinweb}

\begin{grammarlens}[title={What you've built}]
One more word for saying what happens if.
\end{grammarlens}

\subsection*{Your turn}
\begin{itemize}
\raggedright
  \item \emph{Say it:} \emph{vibhinnaḥ}
  \item \emph{Say it:} \emph{vibhinnaḥ}, once more
  \item \emph{Say it:} \emph{yadi matāni vibhinnāni santi, tarhi vayaṃ punaḥ cintayāmaḥ.} — If our opinions differ, we think again.
  \item \emph{Recall:} say the Sanskrit for in two ways, then the Sanskrit for the other one, then say \emph{vibhinnaḥ} again
\end{itemize}

\begin{tcolorbox}[breakable,colback=teal!4,colframe=teal!35!black,title={Before you move on}]
What does \sk{विभिन्नः} mean? (\textquotedblleft{}Different\textquotedblright{}.) Say it once more.
\end{tcolorbox}
\section[(dialogue)]{Looking back: from Coming to Terms to Two Ways to Go}

\label{lesson:SA-R541-review-534-541}



\begin{quote}
Say the Sanskrit for a go-between and an agreement.
\end{quote}

\subsection*{Your turn}
Say these aloud:

\begin{itemize}
\raggedright
  \item a go-between, an agreement, agrees to, turns down, consent
  \item keeps to, true to one's word, trusted, a breaking, one's word
  \item the unseen, what must be, a coincidence, by chance, turns out well
  \item a die, loses at play, a cheat, covetous, a gambler
  \item an omen, destiny, auspiciousness, an astrologer, draws near
  \item heedful, inattention, watchful, suspects, keeps safe
  \item close at hand, left over, completion, a waste of time, spends time
  \item wavers, weighs up, in two ways, the other one, different
\end{itemize}

\begin{tcolorbox}[breakable,colback=teal!4,colframe=teal!35!black,title={Before you move on}]
A go-between? (\textbf{\sk{मध्यस्थः}}, \emph{madhyasthaḥ}.) An agreement? (\textbf{\sk{संविद्}}, \emph{saṃvid}.) Agrees to? (\textbf{\sk{अंगीकरोति}}, \emph{aṅgīkaroti}.) Turns down? (\textbf{\sk{प्रत्याख्याति}}, \emph{pratyākhyāti}.) Consent? (\textbf{\sk{सम्मतिः}}, \emph{sammatiḥ}.) Keeps to? (\textbf{\sk{प्रतिपालयति}}, \emph{pratipālayati}.) True to one's word? (\textbf{\sk{सत्यसन्धः}}, \emph{satyasandhaḥ}.) Trusted? (\textbf{\sk{विश्वस्तः}}, \emph{viśvastaḥ}.) A breaking? (\textbf{\sk{भंगः}}, \emph{bhaṅgaḥ}.) One's word? (\textbf{\sk{वचनम्}}, \emph{vacanam}.) The unseen? (\textbf{\sk{अदृष्टम्}}, \emph{adṛṣṭam}.) What must be? (\textbf{\sk{भवितव्यता}}, \emph{bhavitavyatā}.) A coincidence? (\textbf{\sk{संयोगः}}, \emph{saṃyogaḥ}.) By chance? (\textbf{\sk{यदृच्छया}}, \emph{yadṛcchayā}.) Turns out well? (\textbf{\sk{सम्पद्यते}}, \emph{sampadyate}.) A die? (\textbf{\sk{पाशकः}}, \emph{pāśakaḥ}.) Loses at play? (\textbf{\sk{हारयति}}, \emph{hārayati}.) A cheat? (\textbf{\sk{वञ्चकः}}, \emph{vañcakaḥ}.) Covetous? (\textbf{\sk{लुब्धः}}, \emph{lubdhaḥ}.) A gambler? (\textbf{\sk{द्यूतकारः}}, \emph{dyūtakāraḥ}.) An omen? (\textbf{\sk{शकुनः}}, \emph{śakunaḥ}.) Destiny? (\textbf{\sk{नियतिः}}, \emph{niyatiḥ}.) Auspiciousness? (\textbf{\sk{मंगलम्}}, \emph{maṅgalam}.) An astrologer? (\textbf{\sk{गणकः}}, \emph{gaṇakaḥ}.) Draws near? (\textbf{\sk{उपतिष्ठते}}, \emph{upatiṣṭhate}.) Heedful? (\textbf{\sk{अप्रमत्तः}}, \emph{apramattaḥ}.) Inattention? (\textbf{\sk{अनवधानम्}}, \emph{anavadhānam}.) Watchful? (\textbf{\sk{जागरूकः}}, \emph{jāgarūkaḥ}.) Suspects? (\textbf{\sk{आशंकते}}, \emph{āśaṅkate}.) Keeps safe? (\textbf{\sk{संरक्षति}}, \emph{saṃrakṣati}.) Close at hand? (\textbf{\sk{आसन्नः}}, \emph{āsannaḥ}.) Left over? (\textbf{\sk{अवशिष्टः}}, \emph{avaśiṣṭaḥ}.) Completion? (\textbf{\sk{समाप्तिः}}, \emph{samāptiḥ}.) A waste of time? (\textbf{\sk{कालक्षेपः}}, \emph{kālakṣepaḥ}.) Spends time? (\textbf{\sk{अतिवाहयति}}, \emph{ativāhayati}.) Wavers? (\textbf{\sk{दोलायते}}, \emph{dolāyate}.) Weighs up? (\textbf{\sk{विमृशति}}, \emph{vimṛśati}.) In two ways? (\textbf{\sk{द्विधा}}, \emph{dvidhā}.) The other one? (\textbf{\sk{अपरः}}, \emph{aparaḥ}.) Different? (\textbf{\sk{विभिन्नः}}, \emph{vibhinnaḥ}.)
\end{tcolorbox}
